% Imporditud: tools/import_eksperimendid.py
% Allikas: Eesti füüsikaolümpiaadi eksperimendi ülesannete
%   kogu (Rael Kalda, 2023), ülesanne Ü61, lahendus L61.
\setAuthor{EFO žürii}
\setRound{piirkonnavoor}
\setYear{2013}
\setNumber{G E2}
\setDifficulty{2}
\setTopic{termodünaamika}
\setVahendid{kuum vesi otse veekeedukannust, plasttops, termomeeter, stopper, joonlaud}
\setPunktid{12}
\setAllikas{Eksperimendi kogu Ü61, lahendus lk 59}
\setLahendusseis{toores}

\prob{Vee jahtumine}
Kuidas sõltub topsitäie vee soojuskadude võimsus temperatuurist, kui lasta kuumal veel jahtuda? Joonistage graafik, milles sisaldub temperatuurivahemik 60 $^\circ$C kuni 80 $^\circ$C. Vee erisoojus on 4,2 kJ/(kg$\cdot^\circ$C).
\textit{Vihje:} Lahendades on lubatud lähendada topsi kuju silindriga.

\hint

\solu
Kui kuum vesi on topsi valatud, tuleb sinna asetada termomeeter ning oodata, kuni

selle näit on tõusmise lõpetanud. Vee temperatuuri ajasõltuvuse mõõtmiseks on

täpseim lahendus märkida üles aeg, mis kulub iga täiskraadini jõudmiseks, alus-

tades hiljemalt 81$^\circ$C juurest ning lõpetades kõige varem 59$^\circ$C juures. Sel ajal ei

tohi topsi liigutada ega liigset tuult topsi kohal tekitada. Kindlasti tuleb mõõtmise

ajal topsis olevat vett termomeetriga segada, et selle temperatuur oleks ühtlane

ning termomeeter mõõdaks tõesti kogu vee keskmist temperatuuri.

Teiseks tuleb leida topsis oleva vee mass. Selleks tuleb mõõta veetaseme kõrgus

tkoeps(skims hin, etoppisnidläablaimSõ=õt18p$\pi$õh()dja22s+d1dj12a).ü(lTeävpalnde2v. aLläehmenkdeaskdmesitsoeppsiinsdilainladjraiogak,soonnsSell=e 1 1 $d_1$)2 8 $\pi$ d22 + d12 $-$ 3 ($d_2$ $-$ .) Vee ruumala on siis ligikaudu V = Sh, sellest saame

vee massi m = $\rho$Sh.

Soojuskadude võimsus temperatuuril T on N = Q/$\Delta$t = cm$\Delta$T /$\Delta$t, kus tempera-

tuurivahemiku $\Delta$T keskmine on T ning selle vahemiku läbimiseks kulus aega $\Delta$t.

Piisab, kui N on arvutatud viie erineva T jaoks, näiteks 60$^\circ$C, 65$^\circ$C, 70$^\circ$C, 75$^\circ$C ja

80$^\circ$C. Täpseim on valida $\Delta$T = 2$^\circ$C, ehk siis näiteks 70$^\circ$C-le vastava N leidmiseks

kasutada temperatuurivahemikku 71$^\circ$C - 69$^\circ$C.

Tulemuse graafikul peab soojuskadude võimsus temperatuuri tõustes suurenema.

Arvväärtused sõltuvad kasutatud topsist (ja vee kogusest). Näiteks ühe läbi proovi-

tud topsi puhul tulid soojuskaod 33 W temperatuuril 80$^\circ$C ja 12 W temperatuuril

60$^\circ$C.

Hindamine: Temperatuuri ajasõltuvuse mõõtmise meetodi kirjeldus (1 p). Topsis oleva vee segamine mõõtmise ajal (1 p). 10 või rohkem mõõtepunkti temperatuurivahemikus 58$^\circ$C kuni 82$^\circ$C (2 p). 8..9 mõõtepunkti (1,5 p), 6..7 mõõtepunkti (1 p), 4..5 mõõtepunkti (0,5 p). Topsi keskmise ristlõikepindala leidmine (2 p) (kui on kasutatud ainult põhja või ülemise otsa pindala, siis 1 p ). Vee ruumala arvutamine (1 p). Soojuskadude võimsuse valemi tuletamine (2 p). Soojuskadude võimsuse füüsikaliselt mõistliku graafiku joonistamine (s.h. ühikutega teljed, mõistlik ruumikasutus): graafikul 5 punkti (3 p), 4 punkti (2 p), 3 punkti (1 p).
\probend
